\documentclass[english,course]{lecture}

\usepackage{listings}
\usepackage{cleveref}
%--------------------------------------------------------------------------
% First, provide some data about this document
\title{lecture notes}
\subtitle{Cybersecurity Specialist}
%\shorttitle{Shortened title} % For headers; if undefined, the usual title will be used
\ccode{CS0424IT} % Most of these data are not compulsory
% \subject{Subject of the Talk}
\author{Simone La Porta}
% \spemail{speaker@email.com}
% \email{email@email.com}
% \speaker{Speaker's name}
\date{27}{05}{2024}
\dateend{13}{09}{2024}
% \conference{Lecture hall 7}
% \place{University of Physics}
\flag{\href{https://github.com/simone0509/CS0424IT}{\includegraphics[scale=0.06]{images/github_logo.png}} \\ Written in \LaTeX}
% \attn{Place anything here to gather your readers' attention. This could be a warning, a disclaimer, a license, or, more likely, some helpful suggestions for your readers.}
% \morelink{https://github.com/simone0509/CS0424IT}
%-----------------------------------------------------------------------------


% And then begin your document
\begin{document}

\section{Introduction}
The Cybersecurity Specialist course aims to train professionals in the field of information security. These individuals will possess strong technical skills, providing added value to companies in the fight against cybercrime. The course is divided into three units:
\begin{enumerate}
    \item \textbf{Unit 1} focuses on the theoretical prerequisites and technical skills necessary for an Ethical Hacker. It covers topics such as networking, operating systems, and an introduction to programming;
    \item \textbf{Unit 2} centers on the phases of Penetration Testing, exploring the tools and techniques used by hackers in the real world;
    \item \textbf{Unit 3} provides students with a comprehensive understanding of how to monitor security events, manage ongoing attacks, and adopt best practices at the enterprise level to minimize the impact on business activities.
\end{enumerate}

In the past century, computing and the web were primarily the domain of experts, including hackers. The term \emph{hacker}, often associated with digital piracy, encompasses three distinct types of hackers:
\begin{itemize}
    \item \textbf{White Hat Hackers}, also known as "Ethical Hackers," operate with a strict adherence to ethical standards. Their work involves improving security with the consent of the system owner;
    \item \textbf{Grey Hat Hackers} operate in a legal and ethical gray area. They often act without the owner's permission but with the intent of improving security. While their actions can uncover vulnerabilities, they can also be controversial and sometimes illegal;
    \item \textbf{Black Hat Hackers} are criminals who break into computer networks with malicious intent. They may deploy malware to destroy files, steal information, hold computers hostage, or pilfer passwords, credit card numbers, and other personal data. Their motivations are typically opportunistic, such as financial gain. Stolen data is often sold on the dark web, where items like credit card details, online payment system access, medical records, and even streaming service accounts are traded.
\end{itemize}

An Ethical Hacker is a cybersecurity expert capable of simulating cyber-attacks to identify potential vulnerabilities in a company's systems. These simulations, known as \emph{Penetration Tests}, are crucial for detecting and fixing security issues in digital networks, software, and devices, thereby protecting enterprises and public entities from cybercriminal activities. Key responsibilities of an Ethical Hacker include:
\begin{itemize}
    \item Conducting penetration tests on IT infrastructures and web applications;
    \item Ensuring the security of sensitive and private data, such as payment details, login credentials, and passwords.
\end{itemize}

It is crucial to emphasize that penetration testing should only be performed with the formal consent of the system or network owner. Conducting such tests without permission is illegal and can lead to severe legal consequences.
 % Introduction

\section{Unit 1: Fundamentals of Ethical Hacking}

\subsection{Networking}
\lecture[4 hours]{27}{05}{2024} % Mark new lectures if you like
Pick between \texttt{seminar}, \texttt{talk} or \texttt{course} (this document) to get appropriately different layouts. Seminars being shorter, for example, will not carry the contents section you see above but are otherwise mostly similar to courses. Talks are two-column layouts, which means you cannot have lecture numbers (on your right) or margin notes (see below). Use appropriate options. For more consult \url{http://vhbelvadi.com/latex-lecture-notes-class}.

\subsection{Operating System \& Programming Languages}

\subsection{Python for Hackers and Web Applications}

\subsection{BUILD WEEK 1: Network Security design}
 % Unit 1

And this is a nice \texttt{\$\$...\$\$} display environment:
$$
\Delta v = \int\displaylimits_{t_0}^{t_1} a \, \textrm{d}t
$$
Maecenas ut nisi condimentum nisi iaculis porttitor eu sed metus. Proin faucibus aliquet odio, ac lobortis tortor. Mauris porta molestie tortor blandit pretium. Nulla pulvinar id mauris ut efficitur. Donec posuere tortor a odio pellentesque tincidunt. Nulla mi nunc, accumsan nec lectus ut, euismod vulputate libero.
%
And finally we have the \texttt{align/align*} environment:
\begin{align}
x_f - x_i &= \bar{v}t \nonumber\\
\Rightarrow s &= \bar{v}t
\end{align}
%
\section{Yet another section}
\subsection{And a subsection beneath it}
\lecture[1 hour]{13}{06}{2017}
%
\subsection{And now a subsection}
\subsubsection{With a subsubsection following it}
%
Integer pharetra nulla scelerisque purus luctus iaculis. Mauris pulvinar erat non dui pretium, sed vestibulum sapien condimentum. Nam in urna quis sapien rhoncus placerat vitae sit amet odio. Vivamus finibus euismod nibh vestibulum lobortis.\margintext{These ideas were probably discussed in lecture 1 in a parallel universe.} Integer arcu tortor, vestibulum sit amet iaculis ut, ullamcorper non ante. Pellentesque consectetur nec odio quis placerat. Vestibulum vehicula massa vel euismod blandit.
%
%
\margintext{\protect\vspace{2cm}Table \ref{tab:mori} courtesy of Mori, L.F. `Tables in \LaTeX2$\epsilon$: Packages and Methods'.}
\begin{table}[h]
\centering
\begin{tabular}{clcc}
\toprule
\multicolumn{2}{c}{$D$} & $P_u$ & $\sigma_N$\\
\multicolumn{2}{c}{(in)} & (lbs) & (psi)\\\toprule
\multirow{3}*{5} & test 1 & 285 & 38.00\\\cmidrule(l){2-4}
& test 2 & 287 & 38.27\\\cmidrule(l){2-4}
& test 3 & 230 & 30.67\\\midrule
\multirow{3}*{10} & test 1 & 430 & 28.67\\\cmidrule(l){2-4}
& test 2 & 433 & 28.87\\\cmidrule(l){2-4}
& test 3 & 431 & 28.73\\\bottomrule
\end{tabular}
\caption{A table beautified by the \protect\texttt{booktabs} package.}\label{tab:mori}
\end{table}
%
%
\subsubsection{This subsubsection is all by itself}
%
\vskip7ex
\centering
* * *
%
\end{document}%
